\chapter{L'azienda}
\label{cap:azienda}

\section{Descrizione generale}

\begin{center}
	\includegraphics[height = 4cm]{wavelop-logo.png}
\end{center}

\noindent Wavelop S.R.L.S. è un'azienda nata nel 2018 con sede a Treviso che si occupa di consulenza e di progetti web e mobile per piccole e medie imprese, conta circa 10 dipendenti e offre un ambiente lavorativo giovane e dinamico. Favorisce inoltre lo smartworking e la flessibilità degli orari.\\
Dal 2020 propone progetti per studenti universitari interessati ad un tirocinio formativo che verranno poi  inseriti all'interno di un team di lavoro presso l'azienda.

\section{Modello di svluppo}

Wavelop segue un modello di sviluppo \textit{agile} con metodologia \textit{scrum}. Di seguito le caratteristiche che contraddistinguono questo metodo di lavoro:

\begin{itemize}
  \item il lavoro viene diviso in archi temporali della durata di una settimana, denominati \textit{sprint};
  \item all'inizio di ogni \textit{sprint} si tiene lo \textit{sprint planning}, una riunione in cui si chiarisce quali obiettivi, espressi sotto forma di \textit{user stories}, si mira a portare a termine;
  \item alla conclusione di ogni \textit{sprint} si tiene invece lo \textit{sprint review}, una riunione in cui si discute cosa si è fatto durante lo \textit{sprint}, mostrando i risultati ottenuti, di solito tramite brevi demo;
  \item l'attività giornaliera inizia con la partecipazione allo \textit{stand-up meeting}, una breve riunione della durata di circa 15 minuti, dove, a turno, ogni membro del team espone brevemente cosa ha fatto il giorno precedente e cosa intende fare oggi, oltre che mettere al corrente tutti di eventuali ostacoli incontrati nello sviluppo.
\end{itemize}

\noindent Adottare una metodologia di questo tipo permette una comunicazione costante fra i membri del team e i clienti che porta a un tracciamento chiaro e preciso dell'avanzamento del progetto, mettendo subito in evidenza eventuali incomprensioni, permettendo quindi delle correzioni tempestive, in accordo col cliente.